%=================================================================
% Respuesta a los revisores -- Segundo Corte -- Equipo 0 (EJEMPLO)
% LabDST RoboCoffee, ITESO.
%
% Antes de compilar este documento, compilar el paper revisado:
%   "../Equipo 0 - segundo corte/Equipo 0 - segundo corte.tex"
%   (pdflatex -> bibtex -> pdflatex -> pdflatex)
% Despues compilar ESTE documento dos veces (pdflatex x2).
%
% La carta se escribe en ingles, como el paper. El comentario del
% revisor se copia TEXTUAL desde el PDF de observaciones (en este
% curso las observaciones estan en espanol).
%=================================================================
\documentclass{LabDSTresponse}
\LabDSTpaper{"../Equipo 0 - segundo corte/Equipo 0 - segundo corte"}

\begin{document}
\LabDSTtitle{Second Cut}{Team 0}{A Load-Cell Cup Detection Module}

\LabDSTsection{Acknowledgements}
The authors thank the reviewers for their careful evaluation. Line,
section, and figure numbers refer to the revised manuscript.
\textbf{Main changes:} the Introduction now states a research question
(Section~\ref{sec:intro}), and Section~\ref{sec:methods} adds the
module architecture (Fig.~\ref{fig:arch}).

\LabDSTsummary

%-----------------------------------------------------------------
\LabDSTreviewer{1}{Format, LaTeX and template compliance}

\LabDSTcomment{R1.1}{A}{Section~\ref{sec:intro}, \rln{intro-cite}, \rln{rw-rfid}}
{Las fuentes están escritas a mano en formato autor--año.}
{We created one BibTeX entry per source and replaced every hand-written citation with
\texttt{\textbackslash cite\{key\}} (Section~\ref{sec:intro}, \rln{intro-cite}, \rln{rw-rfid}).}

\LabDSTcomment{R1.2}{A}{Whole manuscript}
{Los párrafos terminan con \texttt{\textbackslash\textbackslash}.}
{All forced line breaks were removed; paragraphs are separated by a blank line.}

%-----------------------------------------------------------------
\LabDSTreviewer{2}{Writing and structure}

\LabDSTcomment{R2.1}{A}{\Rlns{rq}{contrib-end}}
{El artículo no formula una pregunta de investigación ni declara sus contribuciones.}
{We added a research question and three contributions at the end of the Introduction
(\rlns{rq}{contrib-end}). The research question now reads:
\LabDSTquote{Can a single load cell under the cup holder detect cup placement and
removal reliably enough, and fast enough, to release the robot automatically?}}

%-----------------------------------------------------------------
\LabDSTreviewer{3}{Technical content}

\LabDSTcomment{R3.1}{A}{Fig.~\ref{fig:arch}, Eq.~\eqref{eq:rule}, \rlns{mm-arch}{mm-thr-end}}
{No se presenta la arquitectura del módulo ni la regla de detección.}
{The architecture is shown in Fig.~\ref{fig:arch} (\rlns{mm-arch}{mm-arch-end}), and the
detection rule with two thresholds is given in \eqref{eq:rule} (\rlns{mm-thr}{mm-thr-end}).}

\LabDSTcomment{R3.2}{P}{Table~\ref{tab:protocol}, \rlns{res-trials}{res-end}}
{No hay pruebas que respalden la confiabilidad de la detección.}
{The test protocol was added in Table~\ref{tab:protocol}, and the first 20 cycles are
reported in \rlns{res-trials}{res-end}. The complete campaign of 100 cycles will be
reported in the next cut.}

%-----------------------------------------------------------------
\LabDSTreviewer{4}{D-Lab and state of the art}

\LabDSTcomment{R4.1}{A}{\Rlns{rw-bench}{rw-end}}
{No se explica por qué RFID no es adecuado para RoboCoffee.}
{Related Work now explains that RFID needs a tag on every disposable cup
(\rlns{rw-bench}{rw-end}).}

\LabDSTcomment{R4.2}{N}{\Rlns{disc-empty}{disc-end}}
{Implementen una segunda versión del módulo con cámara y comparen ambas.}
{We agree that a camera could identify the cup independently of its weight. However,
image processing and controlled lighting are outside the scope of this module. Instead,
vision-based detection is discussed as future work in \rlns{disc-vision}{disc-end}.}

\end{document}
